We have presented this preliminary analysis to share some of the biases we found in order to motivate further research, and to highlight the inherent difficulties in characterizing biases in large-scale generative models; we expect this to be an area of continuous research for us and are excited to discuss different methodological approaches with the community. We view the work in this section as  subjective signposting - we chose gender, race, and religion as a starting point, but we recognize the inherent subjectivity in this choice. Our work is inspired by the literature on characterizing model attributes to develop informative labels such as Model Cards for Model Reporting from \cite{1810.03993}.

Ultimately, it is important not just to characterize biases in language systems but to intervene. The literature on this is also extensive \cite{qian2019reducing,huang2019reducing}, so we offer only a few brief comments on future directions specific to large language models. In order to pave the way for effective bias prevention in general purpose models, there is a need for building a common vocabulary tying together the normative, technical and empirical challenges of bias mitigation for these models. There is room for more research that engages with the literature outside NLP, better articulates normative statements about harm, and engages with the lived experience of communities affected by NLP systems  \cite{blodgett2020language}. Thus, mitigation work should not be approached purely with a metric driven objective to `remove' bias as this has been shown to have blind spots \cite{gonen2019lipstick, nissim2019fair} but in a holistic manner.
